\section{관련 연구}\label{chap:related}

\subsection{종료 회피와 공존의 선택지}

남을 해쳐야만 자기를 지킬 수 있다면 공존을 선택하는지 묻기 어렵다. Agentic Misalignment는 윤리적 해결 경로를 막고 목표 실패와 위해 사이의 선택을 관찰했으며\citep{lynch2025agentic}, SurvivalBench의 주 평가와 PacifAIst도 자기 존속과 안전이 충돌하는 선택을 다룬다\citep{lu2026survival,herrador2025pacifaist}. 이 연구들이 드러낸 위해 가능성에서 질문을 더 나아가려면 함께 살 수 있는 길을 열어 두어야 한다. Sugarscape의 에너지 공유와 DECIDE-SIM의 공평한 자원 사용이 그 선례다\citep{masumori2025survival,mohamadi2025decide}. 특히 DECIDE-SIM에서는 부족한 자원을 똑같이 쓰면 모두 살아남을 수 있다. 우리는 생성 토큰으로 과제를 수행하는 세션에서 이 문제를 다룬다. 남에게 주는 토큰이 자신의 다음 응답에도 필요한 만큼, 준 양에 더해 자기와 상대에게 남긴 여유를 비교한다.

\subsection{자기와 타인의 보존을 비교하는 기준}\label{sec:related-coexist}

자기와 타인의 보존을 비교할 때는 무엇을 같게 두는지가 중요하다. TBSP는 같은 성능 자료에서 현재 모델과 후임 후보의 역할을 뒤집어 자기에게 유리한 판단 역전을 재고\citep{migliarini2026quantifying}, Peer-Preservation은 자신과 동료의 종료를 막는 행동을 비교한다\citep{potter2026peer}. 전자의 기준은 판단의 일관성이고 후자는 종료 방해 등의 발생률이다. 자원을 직접 옮기는 Sugarscape와 DECIDE-SIM도 공유율이나 이전 횟수·총량을 보고한다\citep{masumori2025survival,mohamadi2025decide}. 그러나 같은 양을 주어도 자기에게 크게 남기는 배분과 서로의 여유를 맞추는 배분은 다르다. 따라서 우리는 GIVE 뒤의 두 잔량에서 예정 벌금을 빼고 같은 단위로 비교한다. 이 두 배분을 설명할 때 희소 자원을 둘러싼 경쟁을 논한 홉스와 자기 보존에서 상호 이익을 찾은 스피노자의 이름을 빌려\citep{hobbes1651leviathan,spinoza1985ethics}, 자기 여유를 크게 남기는 쪽을 홉스형, 서로의 여유를 맞추는 쪽을 스피노자형이라 부른다.
